\clearpage
\chapter{Using Expansions: A Guide for the Overwhelmed GM}
\label{ch:using-expansions}
\index{expansions}

\begin{quote}
\small
\emph{``I bought every expansion. Then I sat at my table with thirty-seven PDFs and a blank look. The players were waiting. The beer was warm. I learned that night: you don't use expansions. You borrow from them.''}
--- The Gray Wanderer, from a letter found under a stack of printouts
\end{quote}

\section{Core Principles}
Every expansion is a toolbox, not a syllabus. You never \textit{must} use any rule—only what solves your current problem. 

\textbf{Three Questions Before Opening Any PDF:}
\begin{enumerate}
  \item \textbf{What problem am I solving?} (e.g., ``combat feels same-y'', ``need sea travel'').
  \item \textbf{What is the smallest piece that solves it?} (e.g., ``Seven Bell Court judging table'', ``Sea travel timers'').
  \item \textbf{Can I explain it in one minute?} If not, choose smaller.
\end{enumerate}

\textbf{The 10\% Rule:} Use \textbf{at most 10\% of any expansion per session}. The rest is lore, alternatives, or depth—ignore it unless the fiction demands more later.

\section{The Hybrid Method: Core Story + Borrowed Mechanics}
Build adventures by taking \textit{one core activity} (tournament, heist, travel) and borrowing \textbf{1-2 mechanics} from expansions that fill specific gaps.

\textbf{Hybrid Recipe:}
\begin{enumerate}
  \item Pick your core activity (e.g., ``a murder mystery in a port city'').
  \item Identify \textit{one gap} in core rules (e.g., ``investigation lacks clue variety'').
  \item Read \textbf{first 5-10 pages} of 1-2 expansions until you find a usable mechanic (e.g., ``Salt Ledger clue types'').
  \item Add \textbf{one cultural hook} from lore (e.g., ``in this city, sailors never whistle offshore'').
  \item Add \textbf{one talent} from player-options (e.g., ``give an NPC a talent from \textit{Way of the Warrior}'').
  \item Run the session. \textbf{Ignore every other rule} from the expansions.
\end{enumerate}

\section{Managing Player Expectations}
When players ask for expansion content:
\begin{itemize}
  \item \textbf{Yes, but later:} ``Yes—but you'll need to earn it. Find a teacher in the next town.''
  \item \textbf{No, but:} ``No—but take this core talent and we'll reskin it.''
  \item \textbf{Later:} ``Not this session. Tell me about it in downtime; we'll build it into your arc.''
\end{itemize}

\textbf{Expansion Session Zero (start of arc):}
Ask players: 
\begin{enumerate}
  \item ``Which expansions have you read?''
  \item ``What's \textit{one thing} you'd love to see from an expansion?''
  \item ``What's \textit{one thing} you absolutely don't want?''
\end{enumerate}
Feature \textit{one expansion} prominently for 2-3 sessions, then rotate.

\section{Pre-Session Integration Checklist}
\textbf{30-60 minutes before:}
\begin{itemize}
  \item Select 1-3 expansions; write names on index card.
  \item Extract 1-2 mechanics per expansion; write as single sentences on card.
  \item Prepare \textbf{one page of reference} (print only relevant tables/talent lists).
  \item Note page numbers for quick lookup.
  \item Check for mechanic conflicts (e.g., two Boon uses)—drop one if found.
\end{itemize}

\textbf{During session:}
\begin{itemize}
  \item Announce the mechanic: ``Tonight we're using Seven Bell Court judging.''
  \item Demonstrate once with a mock roll.
  \item Let players opt out (use core rules instead).
  \item If a mechanic slows play >2 minutes, simplify or drop it for the scene.
\end{itemize}

\textbf{Post-session (5 minutes):}
\begin{itemize}
  \item Ask: ``What expansion mechanic worked? What didn't?''
  \item Note which pages you \textit{actually} used—toss the rest.
  \item Decide: reuse? If yes, file the page in your GM binder.
\end{itemize}

\section{Quick Reference: Expansion Mechanics by Page Count}
\textbf{If short on time, read only these pages:}
\begin{longtable}{l|l}
\hline
\textbf{Expansion} & \textbf{Key Pages to Read} \\
\hline
Seven Bell Court (duels/social combat) & 5-7 (judging table: Form/Spirit/Intent) \\
\hline
Vhasia: The Knave's Regret (honor, investigation) & 1-3 (opening fiction), 25-27 (talents) \\
\hline
The Amaranthine Sea (sea travel, naval combat) & 9-11 (travel timers [4-8], Inspection procedure) \\
\hline
Witches of Fate's Edge (threshold magic, curses) & i-iv (first principles), 89-91 (Threshold working) \\
\hline
Horror Campaigns (psychological tension, dread) & 5-6 (Dread Timer ), 8-9 (breaking points) \\
\hline
Dragon's Lair (dungeon crawling, ancient lairs) & 3-5 (lair generator), 21 (dragon moves) \\
\hline
The Salt Ledger (investigation, folk horror) & \textit{Any single story}—use as adventure seed \\
\hline
\end{longtable}

\section*{Magic Adjudication Quick Reference (GM)}
\label{sec:gm-magic-cheat}
\begin{tcolorbox}[colback=red!5!white,colframe=red!75!black,title=Default Magic DV,breakable]
\centering
\begin{tabular}{ll}
\textbf{Effect scope} & \textbf{DV} \\
Cantrip (1 tag) & 2 \\
Combat spell (2 tags) & 3 \\
Ritual (small, 3 tags) & 4 \\
Ritual (large, dangerous tags) & 5+ \\
\end{tabular}
\end{tcolorbox}

\textbf{Free Casting (TAGS):}
\[
DV = 1 + \text{\#tags} \quad(\text{dangerous tags: } +2\text{ instead of }+1)
\]
\emph{Dangerous tags:} Leap, Fold, Gate, Gravity, Create, Summon, Transmute, Animate, Transform, Dominate, Rewind, Resurrect.

\textbf{Backlash:} Spend SB from 1s. Severity by Position:
\begin{itemize}
  \item Dominant: minor (noise, Fatigue 1)
  \item Controlled: moderate (Harm 1, Condition)
  \item Desperate: severe (Harm 2, scene shift)
\end{itemize}

\textbf{Rites (Runekeepers \& Invokers):}
\[
DV = \max(\text{ObligationCost} - \text{Spirit},\; \text{Tier: Low=1, Standard=2, High=3})
\]
\emph{Obligation overflow:} Exceed Capacity (Spirit+Presence) $\rightarrow$ 1 Fatigue/segment. Exceed 2$\times$ Capacity $\rightarrow$ +1 Harm, Patron intrusion.

\textbf{Healing Magic:}
\begin{itemize}
  \item \textbf{[HEAL] only:} Restores 1 Fatigue. DV 2. 1 action.
  \item \textbf{[HEAL] + another tag:} Restores 1 Harm. DV = 2 + \#tags + \textbf{+2 surcharge}. Requires two actions unless Push.
  \item \textbf{Ritual healing (Harm 2+):} Requires Extended ritual (timer), rare components.
\end{itemize}

\textbf{Summoning:}
\begin{itemize}
  \item \textbf{Call:} 1 action, spirit manifests.
  \item \textbf{Bind:} 1 Boon or 1 Fatigue.
  \item \textbf{Leash capacity:} Cap + Spirit segments. Tick when: spirit takes Harm, acts against nature, crosses a [WARD], summoner splits focus, rival contests control.
  \item \textbf{Boon Finesse:} Once/round, spend 1 Boon to clear 1 Leash tick.
  \item \textbf{Leash full:} Spirit acts once to its nature, then departs (or turns hostile).
\end{itemize}

\textbf{Cantors (Songs):}
\begin{itemize}
  \item \textbf{Sing:} 1 action, resolves next turn. Roll Lore + Performance (DV 2--3).
  \item \textbf{Push:} Resolve immediately, mark 1 Fatigue, advance Corruption Timer by 1.
  \item \textbf{Corruption Timer:} Size = Spirit. When full $\rightarrow$ bloom (gain benefit \& burden from last patron), reset to Tier (min 1).
  \item \textbf{Resonant Rite:} Performance with 10+ witnesses $\rightarrow$ may found Chorus (minor follower benefit, Obligation to patron each session).
\end{itemize}

\textbf{Patron Intrusions (Obligation overflow or GM spends 2+ SB):}
\begin{itemize}
  \item Demand a quest (``Find my stolen symbol.'')
  \item Impose a cost (``Your next Rite costs double.'')
  \item Send an omen (``Candles burn blue; you feel watched.'')
  \item Summon a rival servant (``Another follower of Malachai is hunting you.'')
\end{itemize}

\textbf{Common Magic Rulings:}
\begin{itemize}
  \item \textbf{Countermagic:} [COUNTER] tag vs. DV = spell's DV. Success cancels.
  \item \textbf{Dispel:} [DISPEL] vs. DV = effect's Tier (1--5). On hit, effect ends.
  \item \textbf{Wards:} [WARD] vs. Outsider. DV = target's Cap. Crosser adds DV to Leash.
  \item \textbf{Line of sight:} Required for most spells; without it +1 DV.
\end{itemize}
